% Directly inputtable paper section. Author: ChatGPT.
\section{Integer exponent gcds and finite-support Steinberg chains}
\label{sec:steinberg-integer-finite-chain}

\noindent\textbf{Author: ChatGPT.}
This section supplies two algebraic bridges needed by the valuation-contact
route.  First, it constructs the canonical common-exponent decomposition of
every positive integer and separates its coherent and primitive residual
defects.  Second, it realizes contact surfaces in an actual finite-support
coordinate module and proves the exact weighted cell norm and a conditional
finite-chain estimate.  No form of the $abc$ conjecture and no analytic
filling estimate is assumed.

\subsection{The canonical common-exponent decomposition}

For $n>1$, write
\[
 n=\prod_{p\in S(n)}p^{e_p},
 \qquad e_p>0,
\]
where $S(n)$ is nonempty, and put
\begin{equation}
 g(n)=\gcd_{p\in S(n)}e_p,
 \qquad
 u(n)=\prod_{p\in S(n)}p^{e_p/g(n)}.
 \label{eq:sifc-gcd-base}
\end{equation}
For the unit we adopt the convention
\begin{equation}
 g(1)=1,\qquad u(1)=1.
 \label{eq:sifc-unit-convention}
\end{equation}
The second convention does not turn the gcd of an empty exponent family
into one; it only defines the decomposition at the unit.

\begin{proposition}[Integer exponent-gcd decomposition]
\label{prop:sifc-integer-decomposition}
For every positive integer $n$,
\[
 g(n)>0,\qquad n=u(n)^{g(n)},\qquad u(n)>0,
 \qquad \operatorname{rad}(u(n))=\operatorname{rad}(n).
\]
If $n>1$, then $u(n)$ has the same nonempty prime support as $n$ and
\[
 \gcd_{p\mid u(n)}v_p(u(n))=1.
\]
The last assertion is deliberately restricted to $n>1$.
\end{proposition}

\begin{proof}
The assertions at $n=1$ follow from
\eqref{eq:sifc-unit-convention}.  Suppose $n>1$.  Its finite prime support
is nonempty, so the gcd $g(n)$ of its positive exponents is positive and
divides every $e_p$.  Consequently
\[
 g(n)\frac{e_p}{g(n)}=e_p
\]
for every supported prime, and unique factorization gives
\[
 u(n)^{g(n)}
 =\prod_{p\in S(n)}p^{g(n)(e_p/g(n))}
 =\prod_{p\in S(n)}p^{e_p}=n.
\]
Every quotient $e_p/g(n)$ is positive.  Thus no supported prime disappears
and no new prime appears, proving support and radical equality.  Dividing a
finite nonzero family by its gcd makes the gcd of the quotient family one,
which proves the final assertion.
\end{proof}

Set
\[
 h(n)=\log n,\qquad h_u(n)=\log u(n),\qquad
 r(n)=\log\operatorname{rad}(n),
\]
and define
\begin{equation}
 \nu(n)=(g(n)-1)h_u(n),\qquad
 \sigma(n)=h_u(n)-r(n).
 \label{eq:sifc-thicknesses}
\end{equation}

\begin{proposition}[Coherent--residual defect split]
\label{prop:sifc-defect-split}
For every positive integer $n$,
\begin{equation}
 h(n)=g(n)h_u(n),\qquad
 h(n)-r(n)=\nu(n)+\sigma(n),\qquad
 \nu(n),\sigma(n)\geq0.
 \label{eq:sifc-defect-identity}
\end{equation}
At $n=1$, the three heights $h,h_u,r$ and both thicknesses
$\nu,\sigma$ vanish (while the convention gives $g(1)=1$).
\end{proposition}

\begin{proof}
Proposition~\ref{prop:sifc-integer-decomposition} and the logarithm of a
positive integral power give $h(n)=g(n)h_u(n)$.  Since $g(n)\geq1$ and
$h_u(n)\geq0$, one has $\nu(n)\geq0$.  The equality of radicals and the
elementary inequality $\operatorname{rad}(u(n))\leq u(n)$ give
$\sigma(n)\geq0$.  Finally,
\[
 g(n)h_u(n)-r(n)
 =(g(n)-1)h_u(n)+(h_u(n)-r(n)).
\]
The unit statement follows directly from the convention.
\end{proof}

\subsection{An actual finite-support contact module}

Let
\[
 \mathcal D=\mathbb N\to_0\mathbb Z,
 \qquad
 \mathcal E=(\mathbb N\times\mathbb N)\to_0\mathbb Z,
\]
where $\to_0$ denotes finitely supported functions.  For $X,Y\in\mathcal D$
define
\[
 (X\otimes Y)_{p,q}=X_pY_q,\qquad
 X\wedge Y=X\otimes Y-Y\otimes X,
\]
and set
\[
 \Omega(A,B,C)=(A-C)\wedge(B-C)\in\mathcal E.
\]
This ordered-pair model stores both orientations of each alternating
coordinate.  For effective divisors $A,B,C:\mathbb N\to_0\mathbb N$, a
finite set $S$, and nonnegative weights $w_p$, put
\[
 L_{S,w}(A)=\sum_{p\in S}A_pw_p,
 \qquad
 \|\Omega\|_{S,w}
 =\frac12\sum_{p,q\in S}|\Omega_{p,q}|w_pw_q.
\]

\begin{proposition}[Exact finite-support mixed area]
\label{prop:sifc-finite-mixed-area}
If the supports of $A,B,C$ are pairwise disjoint, then
\begin{equation}
 \|\Omega(A,B,C)\|_{S,w}
 =L_{S,w}(A)L_{S,w}(B)
  +L_{S,w}(B)L_{S,w}(C)
  +L_{S,w}(C)L_{S,w}(A).
 \label{eq:sifc-finite-mixed-area}
\end{equation}
\end{proposition}

\begin{proof}
The three-leg expansion gives
\[
 \Omega_{p,q}
 =A_pB_q-A_qB_p+B_pC_q-B_qC_p+C_pA_q-C_qA_p.
\]
At any coordinate at most one of $A_p,B_p,C_p$ is nonzero.  Applying this
observation at $p$ and at $q$ gives the pointwise identity
\begin{align*}
 |\Omega_{p,q}|={}&A_pB_q+B_pA_q+B_pC_q+C_pB_q\\
                  &+C_pA_q+A_pC_q.
\end{align*}
After multiplication by $w_pw_q$ and finite summation, each directed
rectangle factors.  For example,
\[
 \sum_{p,q\in S}A_pB_qw_pw_q=L_{S,w}(A)L_{S,w}(B).
\]
Every cross-leg rectangle occurs in both orientations; the prefactor
$1/2$ removes this duplication and yields
\eqref{eq:sifc-finite-mixed-area}.
\end{proof}

A canonical positive rational cell is a triple $P=(a,b,c)$ of positive
integers such that
\[
 a+b=c,\qquad \gcd(a,b)=1.
\]
It represents $a/c\in(0,1)$.  The sum relation also gives
$\gcd(a,c)=\gcd(b,c)=1$, so the three valuation divisors have pairwise
disjoint prime supports.  Let $S(P)$ be their finite union and take
$w_p=\log p$.  Unique factorization and
Proposition~\ref{prop:sifc-finite-mixed-area} give
\begin{equation}
 \Phi(P):=\|\Omega(P)\|_{S(P),\log}
 =\log a\log b+\log b\log c+\log c\log a.
 \label{eq:sifc-cell-norm}
\end{equation}

Let $M(P)$ be the one-sided radical polarization of the three leg heights.
Let $V(P)$ be the bilinear contact loss formed from the three coherent
thicknesses $\nu$, and let $R(P)$ be the corresponding loss formed from the
three residual thicknesses $\sigma$.  Polarizing
Proposition~\ref{prop:sifc-defect-split} in
\eqref{eq:sifc-cell-norm} proves the exact identity
\begin{equation}
 2\Phi(P)=M(P)+V(P)+R(P).
 \label{eq:sifc-cell-split}
\end{equation}
No estimate is used in \eqref{eq:sifc-cell-split}.

\subsection{Finite chains with an exact boundary}

Let $I$ be finite, let $P_0$ and $P_j$ be canonical positive rational
cells, and let $n_j\in\mathbb Z$.  Assume the exact equality in the integral
finite-support module
\begin{equation}
 \Omega(P_0)=\sum_{j\in I}n_j\Omega(P_j).
 \label{eq:sifc-exact-boundary}
\end{equation}

\begin{theorem}[Concrete exact-boundary filling estimate]
\label{thm:sifc-exact-boundary}
Under \eqref{eq:sifc-exact-boundary},
\begin{equation}
 \Phi(P_0)\leq\frac12\biggl(
   \sum_{j\in I}|n_j|\bigl(M(P_j)+V(P_j)\bigr)
   +\sum_{j\in I}|n_j|R(P_j)
 \biggr).
 \label{eq:sifc-calibrated-bound}
\end{equation}
\end{theorem}

\begin{proof}
Take the finite union $S$ of every prime support occurring in the target
and the chain.  Evaluate \eqref{eq:sifc-exact-boundary} at each
$(p,q)\in S\times S$, cast to the reals, and multiply by
$\frac12\log p\log q$.  The triangle inequality gives
\[
 \Phi(P_0)\leq\sum_{j\in I}|n_j|\Phi(P_j).
\]
Substitute \eqref{eq:sifc-cell-split} for each cell and separate the two
sums.  This is exactly \eqref{eq:sifc-calibrated-bound}.
\end{proof}

Theorem~\ref{thm:sifc-exact-boundary} is conditional only on the displayed
finite-support equality.  The inductively generated equivalence relation
from the permitted rational five-term moves remains \emph{open and
unformalized}: this section neither defines all generators and closure
operations nor proves that every generated chain satisfies
\eqref{eq:sifc-exact-boundary}.  It also proves neither of the two analytic
bounds required by Gate VF.

\subsection{Lean realization and formal boundary}

The companion Lean module listed in Section~\ref{sec:formal}
formalizes Propositions~\ref{prop:sifc-integer-decomposition} and
\ref{prop:sifc-defect-split}, the actual finite-support types
\texttt{FiniteDivisor} and \texttt{FiniteExteriorSurface},
Proposition~\ref{prop:sifc-finite-mixed-area}, the canonical rational-cell
specialization \eqref{eq:sifc-cell-norm}, the split
\eqref{eq:sifc-cell-split}, and
Theorem~\ref{thm:sifc-exact-boundary}.  Its final namespace theorem
instantiates the earlier abstract finite-coordinate triangle inequality.
The companion axiom audit
uses no project-specific axiom, and direct compilation with warnings treated
as errors succeeds.
